\title{DeepSeekMath: Pushing the Limits of Mathematical Reasoning in Open Language Models}

\begin{abstract}
Mathematical reasoning remains a substantial obstacle for language models, primarily because of its intricate and highly structured characteristics. This study presents DeepSeekMath~7B, an extension of DeepSeek-Coder-Base-v1.5 7B, further pre-trained on a corpus of 120B math-specific tokens obtained from Common Crawl, which also integrates both natural language and programming data. DeepSeekMath~7B attains a remarkable 51.7\% score on the MATH benchmark at the competition level, all without depending on any external tools or ensemble voting methods, thus reaching proficiency comparable to models like Gemini-Ultra and GPT-4. Furthermore, when evaluating self-consistency with 64 samples, DeepSeekMath~7B achieves 60.9\% on MATH. The advanced mathematical reasoning of DeepSeekMath~7B is underpinned by two principal strategies: first, the effective exploitation of publicly accessible web resources using a carefully crafted data curation pipeline; second, the development of Group Relative Policy Optimization (GRPO)—an adaptation of Proximal Policy Optimization (PPO)—which not only strengthens mathematical reasoning performance but also increases PPO’s memory efficiency.
\end{abstract}

\section{Introduction}

Large language models (LLMs) have transformed methodologies for mathematical reasoning within artificial intelligence, driving notable progress on quantitative reasoning benchmarks \citep{MATH} as well as on geometry reasoning assessments \citep{trinh2024solving}. In addition, LLMs have demonstrated substantial utility in aiding humans to tackle intricate mathematical challenges \citep{tao}. Nonetheless, state-of-the-art systems like GPT-4 \citep{gpt4} and Gemini-Ultra \citep{gemini} remain inaccessible to the public, while available open-source alternatives lag substantially in terms of capability and effectiveness.

This work presents DeepSeekMath, a specialized language model targeting mathematical reasoning, which exhibits superior performance compared to existing open-source models and comes close to the benchmark levels set by GPT-4 in academic evaluations. Central to our approach is the construction of the DeepSeekMath Corpus—a large, high-quality pre-training resource containing 120B math-centric tokens. We compile this dataset from Common Crawl (CC) utilizing a classifier based on fastText \citep{joulin2016fasttext}. Initially, the classifier is trained with positive samples from OpenWebMath \citep{openwebmath} and negatives drawn from a broad assortment of non-mathematical web content. In subsequent cycles, the trained classifier is leveraged to identify new positive samples within CC, which are then curated through human annotation to further enhance dataset quality. These newly annotated instances are used to retrain the classifier, yielding iterative improvements. Model evaluation demonstrates the resulting corpus’s robustness, with DeepSeekMath-Base 7B attaining scores of 64.2\% on GSM8K \citep{gsm8k} and 36.2\% on the challenging MATH dataset \citep{MATH}, both figures surpassing Minerva 540B \citep{minerva}. Furthermore, due to the multilingual nature of the DeepSeekMath Corpus, we observe increased accuracy on Chinese-oriented mathematical tests \citep{wei2023cmath,agieval}. We see our methodology for mathematical data collection and processing as a foundation for further advancements in this domain, recognizing considerable opportunities for continued progress.

DeepSeekMath-Base is built upon DeepSeek-Coder-Base-v1.5 7B \citep{deepseek-coder}, as we find that initializing from a code-focused pretrained model leads to better performance than beginning with a standard LLM. Additionally, we discover that mathematical training not only boosts the model’s aptitude in math, but also leads to notable improvements on general benchmarks such as MMLU \citep{mmlu} and BBH \citep{bbh}, suggesting that math-focused training strengthens the model's overall reasoning skills.

Subsequent to this pre-training stage, we conduct mathematical instruction tuning on DeepSeekMath-Base utilizing datasets featuring chain-of-thought \citep{cot}, program-of-thought \citep{pot, pal}, and tool-integrated reasoning \citep{tora} methodologies. The model obtained through this process, DeepSeekMath-Instruct 7B, surpasses other models of the same 7B scale and achieves performance on par with much larger 70B instruction-tuned open-source models.

In addition, we propose Group Relative Policy Optimization (GRPO), a new reinforcement learning (RL) algorithm inspired by Proximal Policy Optimization (PPO) \citep{schulman2017proximal}. Unlike PPO, GRPO eliminates the need for a critic model by leveraging group-based score estimates to compute the baseline, thus markedly reducing computational demands during training. Utilizing only a portion of English instruction tuning data, GRPO delivers notable performance gains over the competitive DeepSeekMath-Instruct, demonstrated by significant increases on in-domain benchmarks (GSM8K: 82.9\% $\rightarrow$ 88.2\%, MATH: 46.8\% $\rightarrow$ 51.7\%) as well as on out-of-domain mathematical datasets (such as CMATH: 84.6\% $\rightarrow$ 88.8\%) throughout the RL phase. Furthermore, we present a comprehensive framework for interpreting multiple methods—Rejection Sampling Fine-Tuning (RFT) \citep{yuan2023scaling}, Direct Preference Optimization (DPO) \citep{dpo}, PPO, and GRPO—revealing that these approaches can all be understood as manifestations of direct or streamlined RL strategies. We supplement our theoretical analysis with a wide range of empirical studies, including comparisons of online versus offline training, evaluation of outcome versus process supervision, and investigations into single-turn versus iterative RL, to thoroughly dissect the key factors underlying this framework. Finally, we elucidate the mechanisms through which our RL approach enhances the capabilities of instruction-tuned models and discuss future avenues for advancing RL in accordance with the unified paradigm presented.

\subsection{Contributions}
Our contribution includes scalable math pre-training, along with the exploration and analysis of reinforcement learning.

\textbf{Math Pre-Training at Scale}

\begin{itemize}[topsep=0pt]
    \item This study demonstrates that the publicly available Common Crawl resource harbors significant mathematical data. Through a carefully crafted data curation process, we assembled the DeepSeekMath Corpus, a large-scale dataset comprising 120B tokens sourced from web pages selected specifically for mathematical relevance. This corpus is nearly sevenfold larger than the math web page collection employed by Minerva \citep{minerva}, and nine times greater than that of the recently introduced OpenWebMath \citep{openwebmath}.

\item DeepSeekMath-Base 7B, our foundational model trained beforehand, matches the performance level of Minerva 540B \citep{minerva}, suggesting that parameter count alone does not determine a model’s mathematical reasoning ability. High performance can also be attained by more compact models if trained on sufficiently high-quality datasets.

\item We present results from our experiments involving math training. Pre-training models on code before initiating math training enhances their performance on mathematical problem solving, regardless of whether tools are utilized. These outcomes provide some insight into the ongoing debate: \textit{does code training enhance reasoning skills?} Our evidence suggests that code training does bolster reasoning, particularly in the context of mathematics.

    \item Although training on arXiv papers is common, especially in many math-related papers, it brings no notable improvements on all mathematical benchmarks adopted in this paper.
\end{itemize}

\textbf{Exploration and Analysis of Reinforcement Learning}

\begin{itemize}[topsep=0pt]
    \item We propose Group Relative Policy Optimization (GRPO), a novel reinforcement learning algorithm that is both resource-efficient and effective. Unlike traditional approaches, GRPO eliminates the need for a critic network by calculating the baseline from group scores, leading to a considerable reduction in training overhead when compared to Proximal Policy Optimization (PPO).
    \item We show that implementing GRPO markedly improves the instruction-tuned performance of DeepSeekMath-Instruct, leveraging only the instruction-tuning dataset. Additionally, we note progress in out-of-domain generalization throughout the reinforcement learning stage.
    \item We present a unified framework for understanding various algorithms, including RFT, DPO, PPO, and GRPO. Furthermore, we conduct a thorough experimental investigation, comparing settings such as online versus offline learning, supervision of outcome versus process, and single-turn against iterative reinforcement learning, in order to analyze the core components of this framework.
    \item Utilizing our unified framework, we analyze why reinforcement learning is effective in this context and outline several promising avenues to further advance reinforcement learning approaches for large language models (LLMs).
\end{itemize}

\subsection{Summary of Evaluations and Metrics}

\begin{itemize}[topsep=0pt]
    \item \textbf{English and Chinese Mathematical Reasoning}: We perform extensive evaluations of our models on both English and Chinese mathematical benchmarks, spanning tasks from elementary to collegiate difficulty. For English, the evaluated benchmarks include GSM8K \citep{gsm8k}, MATH \citep{MATH}, SAT \citep{llemma}, OCW Courses \citep{minerva}, and MMLU-STEM \citep{mmlu}. The Chinese assessments make use of MGSM-zh \citep{mgsm}, CMATH \citep{wei2023cmath}, Gaokao-MathCloze \citep{agieval}, and Gaokao-MathQA \citep{agieval}. Our analysis covers two capabilities: producing comprehensive text-based solutions independently and tackling problems utilizing Python.

On assessments conducted with English benchmarks, DeepSeekMath-Base rivals the performance of the proprietary Minerva 540B model \citep{minerva}, and significantly outperforms all available open-source base models—including Mistral 7B \citep{mistral} and Llemma-34B \citep{llemma}—irrespective of their exposure to math-specific pre-training, often by a wide margin. Remarkably, DeepSeekMath-Base also leads on Chinese benchmarks, a result that can likely be attributed to our inclusion of high-quality mathematical data in various languages, rather than restricting pre-training datasets solely to English as done in previous studies \citep{minerva,llemma}. After applying mathematical instruction fine-tuning and reinforcement learning, both DeepSeekMath-Instruct and DeepSeekMath-RL achieve strong results, surpassing 50\% accuracy on the highly challenging, competition-grade MATH dataset—a milestone reached for the first time by open-source models.

\item \textbf{Formal Mathematics}: We assess DeepSeekMath-Base's abilities on the informal-to-formal theorem proving task introduced in \citep{dsp_proof}, utilizing the miniF2F dataset \citep{minif2f} and employing Isabelle \citep{isabelle} as the proof assistant. DeepSeekMath-Base achieves notable results in few-shot autoformalization scenarios.
\item \textbf{Natural Language Understanding, Reasoning, and Code}: To provide a thorough assessment of the models' overall proficiency in understanding, reasoning, and coding, we test DeepSeekMath-Base on the Massive Multitask Language Understanding (MMLU) benchmark \citep{mmlu}, which covers 57 multiple-choice tasks across varied academic disciplines; the BIG-Bench Hard (BBH) benchmark \citep{bbh}, composed of 23 complex tasks that predominantly necessitate multi-step reasoning; and on HumanEval \citep{codex} as well as MBPP \citep{mbpp}, which are widely recognized datasets for measuring code generation performance. Pre-training with mathematical content enhances both language understanding and reasoning abilities.

\section{Math Pre-Training}

\subsection{Data Collection and Decontamination} This section describes the methodology employed to develop the DeepSeekMath Corpus utilizing data sourced from Common Crawl. As illustrated in Figure \ref{fig:data_collect}, we introduce a multi-step, recursive framework that enables the efficient extraction and assembly of an extensive mathematical dataset from Common Crawl, beginning with an initial seed dataset—typically a concise, high-quality compilation of mathematics resources. Importantly, this pipeline is adaptable for assembling large corpora in other areas, such as programming.

Initially, we select OpenWebMath \citep{openwebmath}—a curated repository of high-quality mathematical web documents—as the foundational seed dataset. From this basis, we train a fastText model \citep{joulin2016fasttext} with the goal of retrieving additional web pages that exhibit characteristics similar to those in OpenWebMath. For model training, 500,000 examples are randomly drawn from the seed dataset to serve as positive samples, while another 500,000 pages from Common Crawl are designated as negatives. We utilize a publicly available implementation\footnote{\url{https://fasttext.cc}} for training, setting the vector size to 256, the learning rate at 0.1, maximum word n-gram size to 3, a minimum word occurrence threshold of 3, and running training for three epochs. To manage the scale of Common Crawl, URL-based deduplication and near-deduplication strategies are applied, reducing it to 40B HTML web pages. The trained fastText model is subsequently employed to identify mathematical content from this deduplicated dataset. We mitigate low-quality mathematical documents by ranking all candidate pages using the fastText prediction scores and retaining only the highest-scoring documents. To determine a suitable data cutoff, we conduct pre-training experiments using the top 40B, 80B, 120B, and 160B tokens. For the initial round, we select the top 40B tokens for further use.

Following the initial round of data acquisition, a significant number of mathematical web pages remain unretrieved, primarily because the fastText model was initially trained on a relatively homogeneous set of positive instances. To address this limitation, we seek out further mathematical web domains to expand the seed dataset, with the intention of enhancing fastText model performance. The entire Common Crawl is first partitioned according to distinct domains, each defined as a group of web pages sharing a common base URL. We then compute, for each domain, the proportion of pages already gathered in the initial round. Any domain for which more than 10\% of its pages were collected is categorized as related to mathematics (e.g., \url{mathoverflow.net}). 
Next, we perform manual annotation to identify specific URLs within these math-associated domains that contain mathematical material (for instance, \url{mathoverflow.net/questions}). Uncollected pages that are connected to these mathematically relevant URLs are subsequently incorporated into the seed set. This strategy allows for the extraction of additional positive samples, which serves to further train the fastText model, equipping it to retrieve a greater quantity of mathematical web content in future rounds. Ultimately, after four cycles of data collection, the process yields 35.5M mathematical web pages encompassing a total of 120B tokens. By the fourth round, we observe that approximately 98\% of the mathematical web pages have already been captured by the conclusion of the third round, prompting us to halt further data acquisition.

In order to prevent benchmark contamination, we adopt the methodology described in \cite{deepseek-coder}, excluding web pages that feature questions or answers sourced from English mathematical benchmarks like GSM8K \citep{gsm8k} and MATH \citep{MATH} as well as Chinese benchmarks such as CMATH \citep{wei2023cmath} and AGIEval \citep{agieval}. Our filtering process operates as follows: any fragment of text that contains a 10-gram sequence matching exactly with any substring from the evaluation benchmarks is excised from our mathematical training dataset. Additionally, for benchmark sequences that are under 10 grams in length but no shorter than 3 grams, we apply a precise string matching strategy to eliminate contaminated web content.

\subsection{Validating the Quality of the DeepSeekMath~Corpus}

We run pre-training experiments to investigate how the DeepSeekMath~Corpus is compared with the recently released math-training corpora:

\begin{itemize}[topsep=0pt]
    \item \textbf{MathPile} \citep{mathpile}: This dataset comprises 8.9B tokens compiled from various sources, including textbooks, Wikipedia, ProofWiki, CommonCrawl, StackExchange, and arXiv, with arXiv making up more than 85\% of the total data;
    \item \textbf{OpenWebMath} \citep{openwebmath}: Sourced by filtering CommonCrawl for mathematics-related content, this collection includes 13.6B tokens;
    \item \textbf{Proof-Pile-2} \citep{llemma}: Assembled from OpenWebMath, AlgebraicStack (which contains 10.3B tokens of mathematical code), and arXiv articles (28.0B tokens), this corpus adopts the token ratio between arXiv, web, and code as 2:4:1, consistent with the setup in \cite{llemma}.
\end{itemize}

\subsubsection{Training Setting}
\label{sec:quality-policy}
We fine-tune a general-purpose pre-trained language model containing 1.3 billion parameters—the same architecture as used in the DeepSeek LLMs \citep{deepseek-llm}, referred to here as DeepSeek-LLM 1.3B—specifically for mathematical tasks. Separate versions of the model are trained on each mathematical dataset, processing 150B tokens per corpus. The entirety of the training process utilizes the streamlined HAI-LLM \citep{haillm} infrastructure. In alignment with the training regimen outlined by the DeepSeek LLMs, we employ the AdamW optimizer \citep{adamW}, setting $\beta_1=0.9$, $\beta_2=0.95$, and $\mathrm{weight\_decay}=0.1$. The learning rate follows a stepped schedule: it ramps up to its maximum after an initial 2,000-step warmup period, is reduced to 31.6\% of this value at 80\% completion of training, and drops to 10.0\% of the peak for the final 10\% of training. The highest learning rate used is 5.3e-4, with each batch consisting of 4 million tokens and a sequence context window of 4K tokens.

\subsubsection{Evaluation Results}

\textbf{The DeepSeekMath~Corpus is of high quality, covers multilingual mathematical content, and is the largest in size.}

\begin{itemize}[topsep=0pt]
    \item \textbf{High-quality}:
    We assess the performance of our models on eight downstream mathematical tasks utilizing few-shot chain-of-thought prompting \cite{cot}. As summarized in Table \ref{tab:corpora_comparison}, models trained with the DeepSeekMath~Corpus consistently surpass counterparts on all benchmarks. Figure \ref{fig:corpora_comparisons} further illustrates that at the 50B token mark (corresponding to a single epoch on Proof-Pile-2), the DeepSeekMath-trained model outperforms models trained with Proof-Pile-2, demonstrating the comparatively higher average data quality of the DeepSeekMath Corpus.
    \item \textbf{Multilingual}:
    Data within the DeepSeekMath~Corpus is sourced from a variety of languages, with English and Chinese forming the principal components. According to Table \ref{tab:corpora_comparison}, incorporating the DeepSeekMath~Corpus into training workflows yields considerable improvements in mathematical reasoning for both English and Chinese. Conversely, other available mathematical datasets—which tend to be dominated by English content—show weaker or even detrimental effects on Chinese mathematical task performance.
    \item \textbf{Large-scale}:
    The volume of the DeepSeekMath~Corpus significantly exceeds that of current mathematical datasets. Figure \ref{fig:corpora_comparisons} reveals that DeepSeek-LLM 1.3B, when trained using DeepSeekMath~Corpus, exhibits a much steeper progression in learning and sustains gains over a longer period. By comparison, existing corpora, being much smaller in size and already cycled through several epochs in training, quickly hit a saturation point in terms of achievable model performance.
\end{itemize}

\subsection{Training and Evaluating DeepSeekMath-Base 7B}

This section presents DeepSeekMath-Base 7B, a foundational model characterized by robust reasoning capacity, particularly in mathematical contexts. The model's initialization is based on DeepSeek-Coder-Base-v1.5 7B \citep{deepseek-coder}, and it is further trained using 500B tokens. The dataset comprises 56\% DeepSeekMath Corpus, 4\% AlgebraicStack, 10\% arXiv papers, 20\% code collected from Github, and 10\% natural language content from Common Crawl in both English and Chinese. The training largely follows the procedure outlined in Section \ref{sec:quality-policy}, with two exceptions: the peak learning rate is adjusted to 4.2e-4 and the batch size is set to 10M tokens.

We present an extensive evaluation of the mathematical proficiency of DeepSeekMath-Base 7B, examining its effectiveness in generating independent mathematical solutions without auxiliary resources, its capability to address mathematical tasks utilizing external tools, and its aptitude in formal theorem proving. In addition to its mathematical functionalities, we offer a broader characterization of the foundational model by analyzing its aptitude in natural language comprehension, logical reasoning, and coding abilities.

\paragraph{Mathematical Problem Solving with Step-by-Step Reasoning}

We assess DeepSeekMath-Base's capability to solve mathematical problems via few-shot chain-of-thought prompting \citep{cot}, utilizing eight different benchmarks in both English and Chinese. These benchmarks include tasks centered on quantitative reasoning, such as GSM8K \citep{gsm8k}, MATH \citep{MATH}, and CMATH \citep{wei2023cmath}, as well as multiple-choice question formats, exemplified by MMLU-STEM \citep{mmlu} and Gaokao-MathQA \citep{agieval}. This evaluation spans a wide array of mathematical domains, ranging from elementary-level to college-level material.

As indicated in Table \ref{tab:base_math_cot}, DeepSeekMath-Base 7B achieves the highest scores on all eight benchmarks when compared to other open-source base models, including both the well-known generalist Mistral 7B \citep{mistral} and the more recent Llemma 34B \citep{llemma}, which was math-pretrained on Proof-Pile-2 \citep{llemma}. In particular, for the challenging MATH dataset, DeepSeekMath-Base outperforms all existing open-source base models by a margin exceeding 10\% absolute, and even achieves better results than Minerva 540B \citep{minerva}, a proprietary base model that is 77 times larger, built atop PaLM \citep{palm}, and fine-tuned on mathematical materials.

\paragraph{Mathematical Problem Solving with Tool Use} We assess the ability of models to perform mathematical reasoning with the assistance of programs on GSM8K and MATH datasets, employing a few-shot program-of-thought prompting approach as described in \citep{pot,pal}. For each mathematical question, models generate a Python script, making use of libraries like \textit{math} and \textit{sympy} to handle complex calculations. The output produced by running the generated code serves as the final solution. According to the results in Table \ref{tab:base_math_pot_proof}, DeepSeekMath-Base 7B surpasses the previous best model, Llemma 34B.

\paragraph{Formal Mathematics}

The automation of formal proofs plays a key role in promoting precision and reliability within mathematical reasoning, while also boosting productivity—a trend that has gained substantial momentum in recent years. In our study, we assess the capabilities of DeepSeekMath-Base 7B on the informal-to-formal proof generation task described in \citep{dsp_proof}, which entails transforming an informal statement, its formal equivalent, and an informal proof into a fully formalized proof. Our experiments utilize the miniF2F benchmark \citep{minif2f}, which features Olympiad-level mathematics problems; for each problem, we use few-shot prompting to produce formal proofs in Isabelle. Adhering to the approach in \cite{dsp_proof}, our methodology employs the model to generate preliminary proof sketches, after which we apply the standard automated prover Sledgehammer \citep{sledgehammer} to complete the proofs by addressing any omitted steps. Table \ref{tab:base_math_pot_proof} highlights the robust performance of DeepSeekMath-Base 7B in the context of automated proof formalization.

\paragraph{Natural Language Understanding, Reasoning, and Code}

We assess the capabilities of the models in natural language understanding using MMLU \citep{mmlu}, reasoning through BBH \citep{bbh}, and code-related tasks with HumanEval \citep{codex} and MBPP \citep{mbpp}. As detailed in Table \ref{tab:understanding_reasoning_code}, DeepSeekMath-Base 7B achieves marked improvements on both MMLU and BBH compared to its predecessor, DeepSeek-Coder-Base-v1.5 \citep{deepseek-coder}, underscoring the beneficial effects of targeted mathematical pretraining on both linguistic comprehension and logical reasoning tasks. Furthermore, by integrating code tokens during continued pretraining, DeepSeekMath-Base 7B is able to sustain the coding performance previously established by DeepSeek-Coder-Base-v1.5 across the two programming evaluation sets. In sum, DeepSeekMath-Base 7B demonstrates substantial gains over the general-purpose Mistral 7B \citep{mistral} across all three evaluation categories involving reasoning and coding.

\section{Supervised Fine-Tuning}

\subsection{SFT Data Curation}

We develop an instruction-tuning dataset for mathematics that includes both English and Chinese problems, spanning multiple mathematical domains and a range of difficulty levels. Each problem is matched with solutions presented in several formats: chain-of-thought (CoT) \citep{cot}, program-of-thought (PoT) \citep{pot,pal}, and reasoning approaches that integrate external tools \citep{tora}. The dataset comprises 776,000 training instances in total.

\begin{itemize}[topsep=0pt]
    \item \textbf{English mathematical datasets}: We provide tool-assisted solutions for problems from GSM8K and MATH, and incorporate a filtered portion of MathInstruct \citep{MathInstruct} alongside the Lila-OOD \citep{lila} training set, both containing problems addressed using either CoT or PoT techniques. This English dataset encompasses a wide range of mathematical domains, including but not limited to algebra, probability, number theory, calculus, and geometry.
    \item \textbf{Chinese mathematical datasets}: We curate a collection of Chinese K-12 mathematics questions across 76 detailed sub-topics, such as linear equations, and supply annotations employing both CoT and tool-augmented reasoning approaches.
\end{itemize}

\subsection{Training and Evaluating DeepSeekMath-Instruct 7B}

This section presents DeepSeekMath-Instruct 7B, which is derived from DeepSeekMath-Base via instruction tuning with mathematics-focused tasks. During training, examples are concatenated in random order up to a context window of 4K tokens. The model is optimized over 500 steps, utilizing a batch size of 256 and maintaining a fixed learning rate of 5e-5.

For assessment, we measure the mathematical capabilities of the models both in settings with and without access to external tools, using four quantitative reasoning benchmarks available in both English and Chinese. Our evaluation compares the performance of our model to other leading contemporary systems:

\begin{itemize}[topsep=0pt]
    \item \textbf{Closed-source models} encompass the following: (1) the GPT series, with GPT-4 \citep{gpt4} and GPT-4 Code Interpreter~\footnote{\url{https://openai.com/blog/chatgpt-plugins\#\#code-interpreter}} recognized as the most advanced, (2) Gemini Ultra and Pro \citep{gemini}, (3) Inflection-2 \citep{inflection-2}, (4) Grok-1~\footnote{\url{https://x.ai/model-card}}, along with several recent contributions from Chinese firms, such as (5) Baichuan-3~\footnote{\url{https://www.baichuan-ai.com}}, and (6) the latest iteration, GLM-4~\footnote{\url{https://open.bigmodel.cn/dev/api\#glm-4}}.

    from the GLM family \citep{glm}.

These models are designed for broad, general-purpose tasks and have typically been subjected to extensive alignment processes.
\item \textbf{Open-source models} encompass the following: generic models such as (1) DeepSeek-LLM-Chat 67B \citep{deepseek-llm}, (2) Qwen 72B \citep{qwen}, (3) SeaLLM-v2 7B \citep{seallm}, and (4) ChatGLM3 6B \citep{chatglm3}. In addition, a range of models demonstrate enhanced mathematical capabilities: (5) InternLM2-Math 20B~\footnote{\url{https://github.com/InternLM/InternLM-Math}}, an extension of InternLM2 further trained on mathematical data and subsequently subjected to instruction tuning; (6) Math-Shepherd-Mistral 7B, which employs PPO training \citep{schulman2017proximal} on top of Mistral 7B \citep{mistral} guided by a process-supervised reward model; (7) the WizardMath suite \citep{wizardmath}, which advances mathematical reasoning performance of Mistral 7B and Llama-2 70B \citep{llama2} by leveraging evolve-instruct—a variant of instruction tuning where instructions are generated by AI—and PPO, with training data drawn chiefly from GSM8K and MATH; (8) MetaMath 70B \citep{metamath}, consisting of Llama-2 70B that has been fine-tuned using an augmented dataset combining GSM8K and MATH; (9) ToRA 34B \cite{tora}, a CodeLlama 34B model further trained to integrate mathematical reasoning with external tools; and (10) MAmmoTH 70B \citep{MathInstruct}, a Llama-2 70B model instruction-tuned using the MathInstruct dataset.

As indicated in Table \ref{tab:sft_rl_math}, when evaluated without tool usage, DeepSeekMath-Instruct 7B exhibits robust capabilities in step-by-step reasoning. In particular, on the competition-level MATH benchmark, our model outperforms all available open-source systems and exceeds most proprietary alternatives (such as Inflection-2 and Gemini Pro) by at least 9\% in absolute terms. This superior performance holds even against much larger models (for instance, Qwen 72B) or those refined with specialized math-oriented reinforcement learning methods (including WizardMath-v1.1 7B). Although DeepSeekMath-Instruct performs comparably to Chinese proprietary models like GLM-4 and Baichuan-3 on the MATH dataset, its results remain below those of GPT-4 and Gemini Ultra.

In the evaluation scenario permitting models to combine natural language reasoning with programmatic tool utilization for solving problems, DeepSeekMath-Instruct 7B attains nearly 60\% accuracy on MATH, outperforming all currently available open-source models. Across additional benchmarks, our model rivals DeepSeek-LLM-Chat 67B—the preceding state-of-the-art—despite being just one-tenth its size.
\section{Reinforcement Learning}

\subsection{Group Relative Policy Optimization} Reinforcement learning (RL) has demonstrated its utility in enhancing the mathematical reasoning skills of LLMs following the Supervised Fine-Tuning (SFT) phase \citep{wang2023math,wizardmath}. Here, we present Group Relative Policy Optimization (GRPO), our streamlined and high-performing RL technique.

\subsubsection{From PPO to GRPO} Proximal Policy Optimization (PPO) \citep{schulman2017proximal} is a prevalent actor-critic reinforcement learning framework, frequently employed for RL-based fine-tuning of LLMs \citep{ouyang2022training}. This approach adapts LLMs by maximizing a surrogate objective defined as follows:

\begin{equation}
\footnotesize
    \mathcal{J}_{PPO}(\theta) = \mathbb{E}{[q \sim P(Q), o \sim \pi_{\theta_{old}}(O|q)]} \frac{1}{|o|} \sum_{t=1}^{|o|} \min \left[ \frac{\pi_\theta(o_{t} | q, o_{<t})}{\pi_{\theta_{old}}(o_{t} | q, o_{<t})} A_{t}, \text{clip} \left( \frac{\pi_\theta(o_{t} | q, o_{<t})}{\pi_{\theta_{old}}(o_{t} | q, o_{<t})}, 1 - \varepsilon, 1 + \varepsilon \right)  A_{t} \right] ,
\end{equation}

Here, $\pi_{\theta}$ denotes the current policy model, while $\pi_{\theta_{old}}$ refers to the previous iteration’s policy. The variables $q$ and $o$ represent questions drawn from the question dataset and outputs sampled from the earlier policy $\pi_{\theta_{old}}$, respectively. The hyper-parameter $\varepsilon$, associated with the clipping mechanism in PPO, helps maintain stability during training. The advantage term $A_t$ is estimated using Generalized Advantage Estimation (GAE) \citep{gae}, which relies on the reward sequence $\{r_{\ge t}\}$ and a value function parameterized as $V_{\psi}$. Consequently, PPO involves the joint training of the value function with the policy, and to prevent excessive reward model optimization, a commonly employed technique is to apply a token-wise KL divergence penalty relative to a reference model in the reward for each token \citep{ouyang2022training}, specifically,

\begin{equation}
    r_{t} = r_\varphi(q, o_{\le t}) - \beta \log\frac{\pi_{\theta}(o_{t}|q, o_{<t})}{\pi_{ref}(o_{t}|q, o_{<t})},
\label{eq:PPO-reward}
\end{equation}

where $r_\varphi$  is the reward model, $\pi_{ref}$ is the reference model, which is usually the initial SFT model, and $\beta$  is the coefficient of the KL penalty.

Since PPO usually relies on a value function that is implemented as an additional model roughly equal in size to the policy network, this approach results in notable overhead in both memory usage and computation. Furthermore, within reinforcement learning workflows, the value function serves as a baseline when computing the advantage, thereby reducing variance. In scenarios involving LLMs, however, the reward model tends to allocate reward scores solely to the final token, complicating the learning of a token-wise value function with reliable accuracy. To overcome these challenges, we introduce Group Relative Policy Optimization (GRPO), illustrated in Figure \ref{fig:grpo}. Unlike PPO, GRPO eliminates the reliance on a separate value function. Instead, it leverages the mean reward obtained from a set of generated outputs—each sampled in response to the same prompt—as the baseline. Specifically, for any given question $q$, GRPO generates a collection of outputs $\{o_1, o_2, \cdots, o_G\}$ from the previous policy $\pi_{\theta_{old}}$, and then refines the policy model by maximizing the following objective:

\begin{equation}
\footnotesize
\begin{split}
    \mathcal{J}_{GRPO}(\theta) &= \mathbb{E}{[q \sim P(Q), \{o_i\}_{i=1}^G \sim \pi_{\theta_{old}}(O|q)]}  \\
    & \frac{1}{G}\sum_{i=1}^G\frac{1}{|o_i|} \sum_{t=1}^{|o_i|} \Bigg[ \min \Bigg( \frac{\pi_\theta(o_{i,t} | q, o_{i,<t})}{\pi_{\theta_{old}}(o_{i,t} | q, o_{i,<t})} \hat{A}_{i,t}, \;\; \text{clip} \Bigg( \frac{\pi_\theta(o_{i,t} | q, o_{i,<t})}{\pi_{\theta_{old}}(o_{i,t} | q, o_{i,<t})}, 1 - \varepsilon, 1 + \varepsilon \Bigg) \hat{A}_{i,t} \Bigg) \\ 
    & \hspace{1.5cm} - \beta \mathbb{D}_{KL}\bigg[\pi_{\theta} \,\big|\big|\, \pi_{ref}\bigg] \Bigg],
\end{split}
\label{eq:GRPO-obj}
\end{equation}

where $\varepsilon$ and $\beta$ serve as hyper-parameters. The term $\hat{A}_{i,t}$ denotes the advantage, which is computed using only the relative rewards within each respective group—a method that will be expounded in subsequent subsections. By evaluating advantages relative to group-based outputs, GRPO naturally leverages the comparative dynamics inherent to reward models, which are predominantly trained on datasets comprising comparisons among outputs for the same question. It is noteworthy that, as opposed to incorporating the KL penalty directly into the reward function, GRPO enforces regularization by including the KL divergence between the trained policy and a reference policy within the loss function itself, thus circumventing any added complexity in the computation of $\hat{A}_{i,t}$. Distinct from the KL penalty mechanism formulated in (\ref{eq:PPO-reward}), the KL divergence in this approach is estimated using the following unbiased estimator \citep{kl_approx}:

\begin{equation}
\small
    \mathbb{D}_{KL}\left[\pi_{\theta} || \pi_{ref}\right] = \frac{\pi_{ref}(o_{i,t}|q,o_{i,<t})}{\pi_{\theta}(o_{i,t}|q,o_{i,<t})}- \log\frac{\pi_{ref}(o_{i,t}|q,o_{i,<t})}{\pi_{\theta}(o_{i,t}|q,o_{i,<t})} - 1,
\end{equation}

which is guaranteed to be positive.

\begin{algorithm}[t]
  \small
  \caption{Iterative Group Relative Policy Optimization}
  \textbf{Input} initial policy model $\pi_{\theta_{\text{init}}}$; reward models $r_\varphi$; task prompts $\mathcal{D}$; 
  hyperparameters $\varepsilon$, $\beta$, $\mu$
  \begin{algorithmic}[1]
    \State Initialize policy model $\pi_\theta$ with $\pi_{\theta_{\text{init}}}$
    \For{iteration = 1, \dots, I}
       \State Set reference model $\pi_{ref}$ to current $\pi_{\theta}$
      \For{step = 1, \dots, M}
      \State Draw a mini-batch $\mathcal{D}_b$ from $\mathcal{D}$
      \State Assign $\pi_{\theta_{old}} \gets \pi_{\theta}$
      \State For each prompt $q$ in $\mathcal{D}_b$, generate $G$ candidate responses $\{o_i\}_{i=1}^G$ from $\pi_{\theta_{old}} (\cdot \mid q)$
      \State For each output $o_i$, evaluate the reward $r_i$ using $r_{\varphi}$
      \State Use group relative advantage estimation to compute $\hat{A}_{i,t}$ for each $t$-th token in $o_i$
      \For{GRPO iteration = 1, \dots, $\mu$}
        \State Optimize $\pi_{\theta}$ via the GRPO loss function (Equation \ref{eq:GC-GRPO})
      \EndFor
    \EndFor
    \State Periodically refine $r_\varphi$ utilizing a replay-based training approach.
    \EndFor 
  \end{algorithmic}
  \textbf{Output} $\pi_\theta$
  \label{alg:iter-grpo}
\end{algorithm}

\subsubsection{Outcome Supervision RL with GRPO} Precisely, for any given question $q$, a set of outputs $\{o_1, o_2, \cdots, o_G\}$ is generated using the previous version of the policy model $\pi_{\theta_{old}}$. Each output is evaluated by a reward model, which assigns a score to each, producing $G$ rewards $\mathbf{r}=\{r_1, r_2, \cdots, r_G\}$. These reward values are then standardized by calculating the z-score within the group: subtracting the group's mean and dividing by the group's standard deviation. With outcome supervision, this normalized reward is given as a terminal reward for each output $o_i$, and it is also utilized as the advantage $\hat{A}_{i, t}$ for all tokens in the sequence, explicitly $\hat{A}_{i, t} = \widetilde{r}_i = \frac{r_i- {\rm mean}(\mathbf{r})}{{\rm std}(\mathbf{r})}$. The policy is subsequently updated to maximize the objective function described by equation (\ref{eq:GRPO-obj}).

\subsubsection{Process Supervision RL with GRPO} Outcome supervision constrains the learning signal to a terminal reward after output completion, which may prove inadequate and inefficient for guiding the policy in intricate mathematical scenarios. Building on the approach of \cite{wang2023math}, we incorporate process supervision, which introduces reward signals at the end of each intermediate reasoning step. Specifically, for a question $q$ and a set of $G$ generated responses $\{o_1, o_2, \cdots, o_G\}$, a process reward model evaluates the sequence by assigning step-level scores, resulting in reward collections: $\mathbf{R} = \{ \{r_1^{index(1)},\cdots,r_1^{index(K_1)}\}, \cdots,  \{r_G^{index(1)},\cdots,r_G^{index(K_G)}\} \}$. Here, $index(j)$ identifies the end token position of the $j$-th reasoning step, and $K_i$ counts the total steps in output $i$. To mitigate scale discrepancies, these rewards are standardized via mean and standard deviation normalization: $\widetilde{r}_i^{index(j)} = \frac{r_i^{index(j)} - {\rm mean(\mathbf{R})}}{{\rm std(\mathbf{R})}}$.

Next, the advantage for each token is determined by summing the normalized rewards of all subsequent (or current) steps: $\hat{A}_{i, t} = \sum_{index(j) \ge t} \widetilde{r}_i^{index(j)}$,
after which the policy is trained to maximize the objective set out in equation (\ref{eq:GRPO-obj}).

\subsubsection{Iterative RL with GRPO} During reinforcement learning, the existing reward model may become inadequate for effectively guiding the evolving policy model. To address this, we investigate the use of iterative RL with GRPO. As outlined in Algorithm \ref{alg:iter-grpo}, iterative GRPO involves regenerating training data for the reward model by leveraging samples obtained from the policy model. The reward model is then updated through continued training, utilizing a replay strategy that maintains 10\% of previously collected data. Subsequently, the policy model is assigned as the new reference model, and its training is further conducted using the updated reward model.

\subsection{Training and Evaluating DeepSeekMath-RL}

Our reinforcement learning (RL) experiments utilize DeepSeekMath-Instruct 7B as the foundation. The RL training dataset comprises chain-of-thought-formatted questions pertaining to GSM8K and MATH drawn from the SFT dataset, yielding approximately 144,000 questions. We omit other SFT questions during RL to specifically assess the effects of RL on benchmarks that have not been present in the RL training set. To prepare the training data for the reward models, we follow the methodology outlined in \citep{wang2023math}. The reward model is initialized from DeepSeekMath-Base 7B and trained using a learning rate of $2 \times 10^{-5}$. For GRPO, we employ a policy model learning rate of $1 \times 10^{-6}$, and set the KL regularization coefficient to 0.04. Each prompt is used to generate $64$ output samples. The maximum sequence length is fixed at 1024, and the training batch size is also set to 1024. The policy model undergoes a single parameter update after each round of exploration. We assess DeepSeekMath-RL 7B on various benchmarks in a manner consistent with the evaluation of DeepSeekMath-Instruct 7B. Within this framework, tasks such as GSM8K and MATH that involve chain-of-thought reasoning are classified as in-domain, while all other benchmarks are considered out-of-domain.

Table \ref{tab:sft_rl_math} presents the results achieved by both open-source and closed-source models using chain-of-thought and tool-integrated approaches across English and Chinese evaluation datasets. Key observations include: 1) With chain-of-thought reasoning, DeepSeekMath-RL 7B achieves accuracy rates of 88.2\% on GSM8K and 51.7\% on MATH, outperforming every other open-source model within the 7B–70B range and even most closed-source competitors. 2) Notably, DeepSeekMath-RL 7B is trained solely using instruction tuning data from GSM8K and MATH, in chain-of-thought format, building upon DeepSeekMath-Instruct 7B. Even with this limited training data, it consistently exceeds the performance of DeepSeekMath-Instruct 7B across all evaluation dimensions, highlighting the utility and strength of reinforcement learning methods.

\section{Discussion}
In this section, we will share our findings in pre-training and RL experiments. 

\subsection{Lessons Learnt in Pre-Training}

We begin by detailing our observations during the pre-training phase. Except where noted, our experimental setup remains consistent with the configuration described in Section \ref{sec:quality-policy}. It should also be emphasized that references to the DeepSeekMath Corpus within this discussion pertain to an 89B-token dataset assembled during the second round of data collection.

\subsubsection{Code Training Benefits Mathematical Reasoning}

An often-cited but largely unproven assumption posits that exposure to code enhances reasoning capabilities. Here, we seek to partially address this conjecture in the context of mathematics, demonstrating that code training bolsters models' mathematical reasoning skills regardless of whether they employ external tools.

To study how code training affects mathematical reasoning, we experimented with the following two-stage training and one-stage training settings:

\textbf{Two-Stage Training}

\begin{itemize}[topsep=0pt]
    \item \textbf{Code Training for 400B Tokens $\rightarrow$ Math Training for 150B Tokens}: Our approach involves pretraining DeepSeek-LLM 1.3B using a dataset composed of 400B tokens from code, after which the model undergoes an additional phase with 150B math-related tokens;
    \item \textbf{General Training for 400B Tokens $\rightarrow$ Math Training for 150B Tokens}: For comparative analysis, we conduct a parallel experiment where the initial 400B tokens are sourced from a broad, general-purpose corpus curated by DeepSeek-AI, rather than code, with the goal of discerning the potential benefits that pretraining on code tokens may confer to the model's mathematical reasoning skills versus a foundation of general data.
\end{itemize}

\textbf{One-Stage Training}

\begin{itemize}[topsep=0pt]
    \item \textbf{Math Pretraining with 150B Tokens}: DeepSeek-LLM 1.3B is pretrained on 150B tokens from mathematical datasets;
    \item \textbf{Joint Training with 400B Code Tokens and 150B Math Tokens}: Sequentially fine-tuning on math data after code data negatively impacts code-related abilities. Therefore, we assess whether simultaneously training on a blended dataset—400B code tokens together with 150B math tokens—in a unified stage can promote stronger mathematical reasoning and also mitigate catastrophic forgetting.
\end{itemize}

\paragraph{Results}
Table \ref{tab:code-math} and Table \ref{tab:code-math-general-eval} demonstrate the downstream performance under different training settings.

Incorporating code training enhances program-aided mathematical reasoning, regardless of whether a two-stage or one-stage training approach is used. Table \ref{tab:code-math} demonstrates that, within the two-stage training framework, solely applying code training substantially improves performance on GSM8K and MATH problems executed in Python. Introducing math training in the subsequent stage further boosts results. Notably, adopting a one-stage training regime, where code and math tokens are interleaved, helps to alleviate the problem of catastrophic forgetting associated with two-stage training. This integrated method not only facilitates improved coding capabilities (Table \ref{tab:code-math-general-eval}), but also strengthens program-aided mathematical reasoning (Table \ref{tab:code-math}).

Training on code further enhances mathematical reasoning abilities even when external tools are not employed. In a two-stage training approach, the preliminary phase focused on code delivers noticeable gains on its own. This phase also facilitates greater efficiency during subsequent math-focused training, culminating in optimal outcomes. By contrast, merging code and mathematics tokens into a single-phase training process diminishes mathematical reasoning performance in the absence of tool use. A possible explanation is that, given its relatively small model size, DeepSeek-LLM 1.3B is not sufficiently equipped to comprehensively learn from both coding and mathematical information at once.

\subsubsection{ArXiv Papers Seem Ineffective in Improving Mathematical Reasoning}

ArXiv papers frequently constitute a part of math pre-training datasets \citep{minerva,gpt-f,llemma,mathpile}. Yet, there has been limited systematic investigation into their specific influence on enhancing mathematical reasoning. Interestingly, and somewhat unexpectedly, our experimental results indicate that incorporating arXiv papers does not significantly bolster mathematical reasoning capabilities. To probe this effect, we evaluated multiple model scales, such as DeepSeek-LLM 1.3B and DeepSeek-Coder-Base-v1.5 7B \citep{deepseek-coder}, across arXiv datasets subjected to various preprocessing strategies:

\begin{itemize}[topsep=0pt]
    \item \textbf{MathPile} \citep{mathpile}: a dataset consisting of 8.9 billion tokens, curated through cleaning processes and heuristic filtering methods; notably, scientific papers from arXiv make up more than 85\% of the content.
    \item \textbf{ArXiv-RedPajama} \citep{redpajama}: includes all arXiv LaTeX documents, but with sections such as preambles, comments, user-defined macros, and bibliographies stripped, resulting in a corpus containing 28.0 billion tokens.
\end{itemize}

In our study, we conduct individual training runs of DeepSeek-LLM 1.3B on 150B tokens and DeepSeek-Coder-Base-v1.5 7B on 40B tokens, utilizing distinct arXiv corpora. Our findings suggest that using arXiv papers exclusively does not enhance mathematical reasoning capabilities. Both models, when exposed solely to arXiv-based datasets, fail to demonstrate significant gains or, in some cases, even exhibit reduced performance across a range of mathematical benchmarks varying in complexity. These include quantitative reasoning tasks such as GSM8K and MATH (see Table \ref{tab:arxiv-cot}), multiple-choice assessments like MMLU-STEM (see Table \ref{tab:arxiv-cot}), as well as formal mathematics benchmarks such as miniF2F (refer to Table \ref{tab:arxiv_atp}).

However, this conclusion has its limitations and should be taken with a grain of salt. We have not yet studied:

\begin{itemize}[topsep=0pt]
    \item How arXiv tokens influence math-focused applications beyond the scope of this study, including the task of translating formal theorems or proofs into informal language;
    \item The influence of arXiv tokens when used in conjunction with other datasets;
    \item If the advantages conferred by incorporating arXiv papers would become evident at greater model sizes.
\end{itemize}

Thus, further exploration is required, which we leave for future studies.

\subsection{Insights of Reinforcement Learning}
\subsubsection{Towards to a Unified Paradigm} This section introduces a cohesive framework to examine a variety of training approaches, including SFT, RFT, DPO, PPO, and GRPO. Additionally, we conduct experimental analyses to identify the influential components within this unified structure. Typically, for a given training algorithm, the gradient of the objective function with respect to the model parameters $\theta$ is formalized as follows:

\begin{equation}
    \nabla_{\theta}\mathcal{J}_{\textcolor{red}{\mathcal{A}}}(\theta) = \mathbb{E}[\underbrace{(q,o) \sim \textcolor{red}{\mathcal{D}}}_{Data \ Source}]\left( \frac{1}{|o|} \sum_{t=1}^{|o|}  \underbrace{GC_{{\mathcal{A}}}(q, o, t, \textcolor{red}{\pi_{{rf}}})}_{Gradient \ Coefficient}  \nabla_{\theta}\log \pi_{\theta}(o_t | q, o_{<t})\right).
\label{eq:objective}
\end{equation}

Three primary elements can be identified: 1) \textit{Data Source $\mathcal{D}$}, responsible for specifying the data used in training; 2) \textit{Reward Function $\pi_{{rf}}$}, which provides the training reward signal; and 3) \textit{Algorithm $\mathcal{A}$}, which leverages both the training data and reward signals to compute the gradient coefficient $GC$, thereby dictating the scale of penalization or reward applied to the data. Within this unified framework, we examine a number of prominent approaches:

\begin{itemize}
    \item  \textbf{Supervised Fine-tuning (SFT)}: In SFT, a pretrained model undergoes additional training on datasets that are curated by humans.
    \item \textbf{Rejection Sampling Fine-tuning (RFT)}: RFT extends the training of the SFT model using outputs generated from SFT responses to specific questions, retaining only those outputs that provide correct answers for further fine-tuning.
    \item \textbf{Direct Preference Optimization (DPO)}: DPO enhances the SFT model by conducting further fine-tuning on sampled, augmented outputs, employing a pairwise DPO loss function to guide learning.
    \item \textbf{Online Rejection Sampling Fine-tuning (Online RFT)}: In contrast to standard RFT, Online RFT begins with the SFT model as its policy and improves it by continuously fine-tuning with outputs that are sampled in real-time from the evolving policy model itself.
    \item \textbf{PPO/GRPO}: PPO/GRPO adopts the SFT model as the initial policy, and subsequently employs reinforcement learning, fine-tuning the model based on outputs drawn from the live policy model.
\end{itemize}

We summarize the components of these methods in Table \ref{tab:data_policy}.
Please refer to Appendix \ref{app:analysis-rl} for a more detailed derivation process.

\paragraph{Observation about Data Source} The data source can be categorized into online and offline sampling. Online sampling refers to collecting training data generated from the real-time exploration of the training policy model. In contrast, offline sampling uses data sampled from the outputs of the original SFT model. Approaches such as RFT and DPO utilize offline data, whereas Online RFT and GRPO make use of data obtained through online sampling.

As illustrated in Figure \ref{fig:rl-analysis}, our results indicate that Online RFT achieves notably better performance than RFT across both benchmarks. In the initial training phases, Online RFT and RFT deliver similar outcomes, but as training progresses, Online RFT establishes a clear lead, thereby underscoring the benefits of online training. This observation aligns with expectations: early in training, both the actor and SFT model behave similarly and produce samples with minimal variation. As training advances, samples generated by the actor begin to diverge substantially, making the advantages of real-time data sampling increasingly pronounced.

\paragraph{Observation about Gradient Coefficient} The algorithm leverages the reward signal in determining the gradient coefficient, which in turn is used for updating model parameters. In our experimental framework, we differentiate the reward function into `Rule' and `Model' components. `Rule' assesses response quality strictly by evaluating answer correctness, while `Model' involves training a reward model that assigns scores to responses. The dataset for training this reward model is derived from rule-based judgments. As detailed in Equations \ref{eq:GC-RFT} and \ref{eq:GC-GRPO}, a principal distinction emerges between GRPO and Online RFT: GRPO specifically calibrates the gradient coefficient in response to the magnitude of rewards output by the reward model. This capability enables tailored reinforcement and penalization relative to the extent of each response’s reward value. Conversely, Online RFT does not incorporate such adjustment; it uniformly rewards correct responses and applies no penalties to incorrect ones, treating all correct outputs equivalently regardless of their specific qualities.

Figure \ref{fig:rl-analysis} illustrates that GRPO outperforms online RFT, underscoring the effectiveness of modifying the coefficients for positive and negative gradients. Moreover, GRPO+PS achieves better results than GRPO+OS, suggesting that employing fine-grained, step-sensitive gradient coefficients offers advantages. Additionally, we investigate the impact of iterative RL by performing two iterations in our experimental setup. The results in Figure \ref{fig:iter-rl} reveal that iterative RL markedly enhances performance, with the most substantial improvement occurring during the initial iteration.

\subsubsection{Why RL Works?} In this study, we utilize reinforcement learning on a selected portion of instruction tuning data, observing considerable improvement over models solely trained via instruction tuning. To investigate the underlying reasons for RL’s effectiveness, we analyze the Pass@K and Maj@K accuracy scores for both Instruct and RL-based models across two benchmark datasets. As illustrated in Figure \ref{fig:maj-pass}, reinforcement learning leads to a notable increase in Maj@K accuracy, while Pass@K remains relatively unchanged. These results suggest that the performance gains from RL are mainly due to a more resilient output distribution; put differently, \textbf{the advancement appears to result from elevating the ranking of correct answers within the TopK outputs, rather than from strengthening the model's core capabilities.} In parallel, \citep{wang2023making} have recognized a \textbf{misalignment problem} present in SFT models on reasoning benchmarks, noting that adopting certain preference alignment techniques \citep{yuan2023rrhf,song2023preference,wang2023making} can help SFT models better align with reasoning objectives.

\subsubsection{How to Achieve More Effective RL?}
\label{sec:effec_rl}
Our findings indicate that RL is quite effective for mathematical reasoning applications. Additionally, we introduce a comprehensive framework that offers a cohesive understanding of a range of well-known training approaches. In this framework, each approach is treated as a variant or direct implementation of RL methodology. As outlined in Equation \ref{eq:objective}, there are three principal components at play: Data Source, Algorithm, and Reward Function. We outline several prospective avenues for advancing each of these components in the future.

\paragraph{Data Source} The foundation of any training methodology lies in its data source. Within RL, we define the data source as the collection of unlabeled questions accompanied by their corresponding outputs, which are generated through sampling from the policy model. In our study, the questions utilized are restricted to those drawn from the instruction tuning phase, and response generation employs simple nucleus sampling. We hypothesize that this design choice may contribute to the observation that our RL pipeline yields improvements primarily in Maj@K metrics. Future work will involve extending our RL framework to address question prompts beyond the initial distribution, while incorporating \textbf{advanced sampling (decoding) strategies}, such as tree-search-based techniques \citep{yao2023tree}. Furthermore, \textbf{efficient inference techniques} \citep{xia-etal-2023-speculative,leviathan2023fast,kwon2023efficient,xia2024unlocking}, which critically influence the exploration dynamics of policy models, will also be investigated for their impactful role.

\paragraph{Algorithms} Algorithms utilize both the input data and the reward signal to compute gradient coefficients, which are then used to adjust the model’s parameters. As described in Equation \ref{eq:objective}, contemporary approaches predominantly place full \textbf{TRUST} in the reward function’s signal when modifying the conditional probability assigned to specific tokens. Nonetheless, this reliance presumes the reward signal's accuracy, which cannot be guaranteed—particularly in highly intricate tasks. For instance, even rigorously labeled resources such as the PRM800K datasets \citep{lightman2023let}—annotated by expert human coders—are estimated to contain around 20\% erroneous annotations\footnote{\url{https://github.com/openai/prm800k/issues/12\#issuecomment-1728491852}}. Consequently, there is a need to investigate reinforcement learning techniques that maintain robustness in the presence of noisy or unreliable reward signals. We argue that \textbf{WEAK-TO-STRONG} alignment approaches \citep{burns2023weak} may introduce pivotal innovations to current learning paradigms.

\paragraph{Reward Function} The reward function serves as the primary source of supervision during training. In RL, this function is commonly implemented through a neural reward model. We identify three critical research avenues for reward models: 1) \textbf{Strategies for improving the reward model’s capacity to generalize.} To foster true advancements in LLM capabilities—rather than merely reinforcing existing distributions—the reward model should robustly generalize to unseen questions and sophisticated outputs; 2) \textbf{Mechanisms for representing uncertainty in the reward model.} Incorporating uncertainty has the potential to connect weaker reward models with progression-based learning algorithms, facilitating weak-to-strong learning dynamics; 3) \textbf{Approaches for constructing high-quality, efficient process reward models} that deliver nuanced training feedback tailored to complex reasoning tasks \citep{lightman2023let,wang2023math}.

\section{Conclusion, Limitation, and Future Work}

In this work, we introduce DeepSeekMath, a model that surpasses all available open-source counterparts on the MATH benchmark at the competition level and nears the performance of proprietary systems. DeepSeekMath is based on DeepSeek-Coder-v1.5 7B, followed by continual training using 500B tokens, a considerable portion of which—specifically 120B math tokens—are derived from Common Crawl. Our comprehensive ablation study highlights the untapped potential of web pages as sources of high-quality mathematical data, in contrast to arXiv, which was found to be less advantageous than previously assumed. Furthermore, we present Group Relative Policy Optimization (GRPO), an adaptation of Proximal Policy Optimization (PPO), that enables notable improvements in mathematical reasoning while reducing memory requirements. Experimental results indicate that GRPO maintains its effectiveness even when DeepSeekMath-Instruct 7B achieves top benchmark scores. Additionally, we propose a cohesive framework for interpreting various approaches and identify several promising avenues for advancing reinforcement learning methodologies.

While DeepSeekMath~demonstrates strong performance on benchmarks assessing quantitative reasoning, its abilities in geometry and theorem-proving lag behind those of closed-source counterparts. For example, preliminary testing revealed the model struggles with questions involving triangles and ellipses, suggesting potential biases introduced during both pre-training and fine-tuning data curation. Additionally, due to limitations in model size, DeepSeekMath~falls short of GPT-4 in few-shot learning scenarios. Unlike GPT-4, which exhibits performance gains when provided with few-shot inputs, DeepSeekMath~displays minimal differences between zero-shot and few-shot results. Looking ahead, our objective is to enhance our data selection mechanisms to assemble a more robust and high-quality pre-training dataset. Moreover, we aim to investigate avenues (see Section \ref{sec:effec_rl}) for optimizing reinforcement learning techniques in LLM development.

\end{document}